\documentclass[sigconf,natbib=false]{acmart}

\usepackage{booktabs}
\setlength{\headheight}{15.5pt}
\setlength{\emergencystretch}{3em}
\usepackage{listings}
\usepackage{xcolor}
\usepackage{hyperref}
\usepackage{microtype}

\lstset{
  basicstyle=\ttfamily\small,
  breaklines=true,
  frame=single,
  backgroundcolor=\color{gray!10}
}

%%% arXiv preprint mode — suppress all ACM copyright/conference output
\setcopyright{none}
\settopmatter{printacmref=false, printccs=true, printfolios=true}
\renewcommand\footnotetextcopyrightpermission[1]{}
\acmDOI{}
\acmISBN{}
\acmConference[LICITRA-MMR]{}{}{}
\acmYear{}
\copyrightyear{}

\begin{document}
\sloppy

\title{LICITRA-MMR: A Merkle Mountain Range Ledger Primitive for\\
Cryptographic Runtime Accountability in Agentic AI Systems}

\author{Narendra Kumar Nutalapati}
\affiliation{%
  \institution{Independent Researcher, LICITRA Research Program}
  \city{Portland}
  \state{Oregon}
  \country{USA}
}
\email{narendrakumar.nutalapati@gmail.com}

\begin{abstract}
Agentic AI systems that make autonomous decisions in production create a
new class of accountability requirement: the ability to reconstruct, after
the fact, the exact decision chain with cryptographic verifiability --- a
stronger guarantee than operational logs alone provide. Existing logging
approaches --- including mutable database logs, hash chains, and
general-purpose tamper-evident stores not tailored for agentic governance
--- provide observability but do not explicitly address independently
verifiable runtime decision chains. We present LICITRA-MMR, an open-source
ledger primitive that combines a Merkle Mountain Range (MMR) data structure
with per-organization epoch anchoring, a versioned canonical JSON
specification, and an atomic two-phase commit pipeline. At a block size of
1,000 events, LICITRA-MMR produces inclusion proofs requiring 14 SHA-256
operations, verifies a full epoch chain of 1,000 epochs in under 1~ms, and
supports full leaf rehash of 1 million events in approximately 2 seconds.
The system is fully open-source under MIT license with a tagged release,
reproducible test suite, and evidence bundles. LICITRA-MMR is designed as
a foundational cryptographic audit layer for agentic AI governance systems.
\end{abstract}

\keywords{Merkle Mountain Range, agentic AI, cryptographic audit, runtime
integrity, tamper-evident logging, AI governance}


\begin{CCSXML}
<ccs2012>
<concept>
<concept_id>10002978.10002986.10002990</concept_id>
<concept_desc>Security and privacy~Digital signatures</concept_desc>
<concept_significance>500</concept_significance>
</concept>
<concept>
<concept_id>10002978.10002986.10003006</concept_id>
<concept_desc>Security and privacy~Tamper-proof and tamper-evident designs</concept_desc>
<concept_significance>500</concept_significance>
</concept>
<concept>
<concept_id>10010520.10010521.10010537</concept_id>
<concept_desc>Computer systems organization~Dependable and fault-tolerant systems and networks</concept_desc>
<concept_significance>300</concept_significance>
</concept>
</ccs2012>
\end{CCSXML}

\ccsdesc[500]{Security and privacy~Digital signatures}
\ccsdesc[500]{Security and privacy~Tamper-proof and tamper-evident designs}
\ccsdesc[300]{Computer systems organization~Dependable and fault-tolerant systems and networks}

\maketitle
\fancyhead{}
\fancyhead[LE,RO]{\thepage}
\fancyhead[RE]{\itshape Narendra Kumar Nutalapati}
\fancyhead[LO]{\itshape LICITRA-MMR}
\fancyfoot{}


%% ─────────────────────────────────────────────────────────────────────────────
\section{Introduction}

Regulatory frameworks for high-risk AI systems are converging on a common
implicit technical requirement: that organizations can reconstruct AI
decision chains after an incident with sufficient integrity to satisfy
legal and audit scrutiny. The EU AI Act Article~12~\cite{euaiact} requires
that high-risk AI systems generate logs that are traceable and sufficient
for post-market monitoring. The Colorado AI Act~\cite{coloradoai} (SB~24-205,
effective June~30, 2026) requires deployers to document AI decision-making
processes and conduct impact assessments. Both frameworks are consistent
with an audit capability that most production agentic AI stacks do not
currently provide.

The gap is primarily technical rather than conceptual. Most teams deploy
agentic AI with mutable database logs, cloud logging pipelines, or
dashboard-based observability. These systems answer the question: \textit{what
does our log say happened?} They do not answer: \textit{can we prove the log
was not modified after the incident?} The distinction matters because
administrators can rewrite database rows, cloud logs can be truncated or
filtered, and cleanup scripts can alter historical records. A post-incident
forensic investigation that relies entirely on operator-controlled logs
provides no independent verification guarantee.

Agentic systems intensify this problem. Unlike single-inference ML systems,
autonomous agents make multi-step decisions --- each step potentially
triggering real-world actions --- without per-action human oversight. The
decision chain that produced a harmful outcome may span dozens of sequential
operations, each dependent on prior context. Reconstructing that chain
requires not just that each event was logged, but that the sequence is
complete, ordered, and unmodified.

We make the following contributions:

\begin{enumerate}
\item \textbf{Canonical JSON Specification.} A versioned, cross-implementation
  deterministic serialization format (CANONICAL\_JSON\_SPEC v1.0) that
  ensures identical byte sequences from identical logical payloads across
  independent implementations --- a prerequisite for deterministic verification.

\item \textbf{MMR-Based Per-Organization Tamper-Evident Ledger.} A Merkle
  Mountain Range data structure providing $O(\log n)$ inclusion proofs and
  append-only semantics, with per-organization epoch anchoring that chains
  epoch roots into a hash-linked sequence resistant to retroactive modification.

\item \textbf{Atomic Two-Phase Commit Pipeline.} A propose~$\to$~policy~$\to$
  commit pipeline in which every MMR append, event record, and epoch
  finalization occurs within a single database transaction, preventing partial
  writes that could leave the ledger in an inconsistent state.

\item \textbf{Open-Source Implementation with Reproducible Evidence.} A fully
  tested (11/11 passing), MIT-licensed implementation at tagged release
  v0.1.0-mvp, with reproducible JSON evidence bundles covering five
  tamper-detection scenarios.
\end{enumerate}

%% ─────────────────────────────────────────────────────────────────────────────
\section{Background}

\subsection{Merkle Mountain Ranges}

A Merkle Mountain Range is an append-only authenticated data structure
introduced by Peter Todd in the context of Bitcoin timestamping~\cite{todd2013}.
Unlike many Merkle tree implementations, in which insertion requires updating
internal nodes along the rightmost path, an MMR avoids modifying previously
committed subtrees. It maintains a forest of perfect binary trees (peaks)
that are merged as new leaves are appended. This property makes MMRs ideal
for append-only audit logs: each new element extends the structure in
$O(\log n)$ time and space without modifying any existing node.

Formally, an MMR of size $n$ stores $2n - \text{popcount}(n)$ nodes, where
$\text{popcount}(n)$ is the number of set bits in the binary representation
of $n$. The peaks of the MMR are the roots of the constituent perfect binary
trees. A single MMR root is produced by bagging the peaks: hashing them
right-to-left into a single digest. An inclusion proof for a leaf at
position $p$ consists of the sibling hashes on the path from $p$ to its
subtree peak, plus the remaining peaks, enabling logarithmic verification
without access to the full node set.

At \textsc{block\_size}~$= 1{,}000$ (our default epoch size), the MMR
contains 1,994 nodes and produces inclusion proofs requiring exactly 14
hash operations. This compares favorably to a standard Merkle tree over
the same leaf count, which requires $\lceil\log_2(1000)\rceil = 10$ hashes
per proof path but requires updating internal nodes on insertion. The
formal optimality of the MMR for append-only accumulators was established
by Bonneau et al.~\cite{bonneau2025}.

\subsection{Epoch Anchoring}

LICITRA-MMR partitions the event stream into fixed-size epochs of
\textsc{block\_size} events. When an epoch reaches capacity, an epoch
record is finalized containing:

\begin{lstlisting}
epoch_hash = SHA-256(
  prev_epoch_hash
  || mmr_root
  || SHA-256(metadata))
\end{lstlisting}

where \texttt{prev\_epoch\_hash} is the hash of the preceding epoch (or a
32-byte genesis value for epoch~0), \texttt{mmr\_root} is the bagged peak
hash of the epoch's MMR, and \texttt{metadata} is a canonical JSON document
containing the organization ID, epoch ID, sequence range, event count, and
ledger version.

This construction chains epochs into a hash-linked sequence analogous to a
blockchain header chain, but without distributed consensus overhead.
Modifying any event within an epoch invalidates that epoch's MMR root,
which invalidates all subsequent epoch hashes. An auditor who holds the
current epoch hash and any historical event can verify the complete chain
of custody from genesis to present.

\subsection{The Agentic AI Accountability Gap}

Contemporary agentic AI systems --- systems that autonomously plan and
execute multi-step tasks using large language models as reasoning engines
--- introduce accountability challenges not present in traditional ML
deployments. In a single-inference deployment, a harmful output is a
discrete event that can be logged with a timestamp and input/output pair.
In an agentic system, a harmful outcome may result from a sequence of
individually benign decisions that collectively produce an unintended
effect. Attributing the outcome requires reconstructing the full decision
chain: which prompts the agent saw, which tools it invoked, which
intermediate results it relied upon, and which final action it executed.

Existing logging infrastructure is not designed for this requirement.
Application logs capture events as they occur but provide no
tamper-evidence guarantees. Observability platforms aggregate telemetry
but rely on operator-controlled infrastructure. Neither provides the
forensic integrity consistent with emerging AI regulation: the ability for
an external party to verify, independently of the operator, that a complete
and unmodified record of the decision chain exists.

%% ─────────────────────────────────────────────────────────────────────────────
\section{Threat Model}

LICITRA-MMR is designed to protect the integrity of an append-only audit
ledger against post-hoc modification. We consider the following attacker
classes and explicitly state what the system does and does not protect
against.

\subsection{Protected Against}

\textbf{T1 --- Insider / Database Administrator.} An attacker with direct
write access to the database attempts to modify the \texttt{canonical\_json}
field of a committed event. The leaf hash stored separately in the
\texttt{leaf\_hash} column will no longer match the recomputed SHA-256 of
the modified canonical JSON, and MMR verification will fail at the affected
leaf position.

\textbf{T2 --- Log Scrubber / Row Deletion.} An attacker deletes one or
more event rows to suppress evidence of a specific decision. The monotonic
per-organization sequence number creates a detectable gap: verification
will identify missing sequence positions, and the MMR node set will contain
orphaned nodes at the deleted leaf positions.

\textbf{T3 --- Post-Incident Epoch Modification.} An attacker modifies a
historical epoch hash to hide evidence of prior tampering. Because each
epoch hash is computed over the previous epoch hash, any modification to a
historical epoch invalidates all subsequent epoch hashes.

\textbf{T4 --- Selective History Presentation.} An operator presents only
a subset of the audit log to an external auditor, claiming it is complete.
An independent auditor can verify completeness by recomputing the MMR root
from the presented events and checking it against the claimed epoch hash.
Any omitted events will produce a root mismatch.

\subsection{Not Protected Against}

\textbf{N1 --- Compromised Application Layer.} LICITRA-MMR commits whatever
the application layer submits. Pre-execution integrity is the responsibility
of LICITRA-SENTRY, discussed in Section~\ref{sec:conclusion}.

\textbf{N2 --- Epoch Hash Without Signing.} The current implementation
computes epoch hashes but does not sign them with a private key. Ed25519
epoch signing is planned for v0.2 and discussed in Section~\ref{sec:limitations}.

\textbf{N3 --- Availability Attacks.} LICITRA-MMR does not protect against
denial-of-service attacks against the underlying database or API layer.

\subsection{Trust Assumptions}

The system assumes: (1) SHA-256 is collision-resistant; (2) the database
storage layer preserves byte-level fidelity of stored values; (3) at least
one honest party retains a copy of a recent epoch hash against which
verification can be anchored. This assumption is structurally equivalent
to Certificate Transparency's reliance on at least one honest log
monitor~\cite{laurie2013}. No assumption is made about the trustworthiness
of the operator beyond assumption~(2).


%% ─────────────────────────────────────────────────────────────────────────────
\section{Security Properties}
\label{sec:security}

We formalize the core integrity guarantees provided by LICITRA-MMR.
These are not full proofs; they are precise statements of what the
system guarantees under standard cryptographic assumptions.

\subsection{Append-Only Integrity}

Let $L = (e_1, e_2, \ldots, e_n)$ be the ordered sequence of committed
events for an organization within an epoch. Define the epoch hash as:

\begin{align*}
H_{\text{epoch}} &= \text{SHA-256}(\texttt{prev\_epoch\_hash} \\
  &\quad\|\; \text{MMR\_root}(L) \\
  &\quad\|\; \text{SHA-256}(\texttt{metadata}))
\end{align*}

Under the assumption that SHA-256 is collision-resistant, any adversary
who modifies, deletes, reorders, or inserts an event in $L$ after
$H_{\text{epoch}}$ has been externally anchored must produce either a
second preimage for SHA-256 or a collision in the MMR construction.
Therefore, retroactive modification of committed history without
detection is computationally infeasible under standard hash function
assumptions.

\subsection{Inclusion Soundness}

For any event $e_i$ committed in epoch $k$, the inclusion proof
$\pi_i$ produced by the MMR satisfies:

\[
\text{Verify}(e_i,\; \pi_i,\; \text{MMR\_root}_k) = \texttt{true}
\]

if and only if $e_i$ was included in the committed set $L_k$ at the
time of epoch finalization. Any modification to $e_i$ or to the node
path represented in $\pi_i$ requires either: (1) a collision in
SHA-256, or (2) modification of the persisted node set — both of which
are detectable under the append-only integrity model above.

\subsection{Trust Model Limitation}

These guarantees are relative. If an adversary controls both the
database and all external epoch anchors, they can reconstruct a
consistent but false history. LICITRA-MMR therefore provides
tamper-evidence relative to at least one honest anchor holder,
consistent with the trust assumption stated in Section~3.3.
Epoch signing (planned for v0.2) will strengthen this guarantee
by binding each epoch to a key holder independent of the database.

%% ─────────────────────────────────────────────────────────────────────────────
\section{System Design}

\subsection{Canonical JSON Specification}

A prerequisite for deterministic verification across independent
implementations is that identical logical payloads produce identical byte
sequences. LICITRA-MMR defines a versioned canonical JSON specification
(CANONICAL\_JSON\_SPEC v1.0) with the rules shown in Table~\ref{tab:canon}.

\begin{table}[ht]
\caption{Canonical JSON Specification v1.0}
\label{tab:canon}
\begin{tabular}{lp{4.5cm}}
\toprule
Property & Value \\
\midrule
Character encoding & UTF-8, no BOM \\
Key ordering & Lexicographic by Unicode code point, recursive \\
Key-value separator & \texttt{:} (no spaces) \\
Item separator & \texttt{,} (no spaces) \\
Non-ASCII characters & Preserved as literal UTF-8 \\
NaN / Infinity & Prohibited; raises \texttt{ValueError} \\
Whitespace & None outside string values \\
\bottomrule
\end{tabular}
\end{table}

\textbf{Known limitation.} The specification does not currently normalize
floating-point values. The integer \texttt{1} and the float \texttt{1.0}
produce different byte sequences and therefore different leaf hashes.
Callers must ensure consistent numeric types. Full float normalization per
RFC~8785~\S3.2.2 is planned for v1.1.

\subsection{MMR Structure and Leaf Hashing}

Each event payload is serialized to canonical JSON bytes and hashed:
\[
\texttt{leaf\_hash} = \text{SHA-256}(\text{canonical\_bytes}(\text{payload}))
\]
The leaf hash is appended to the epoch's MMR via \texttt{append\_leaf()},
which adds the leaf node and merges sibling pairs upward. Each merge
produces one parent node:
\[
\texttt{parent\_hash} = \text{SHA-256}(\texttt{left} \| \texttt{right})
\]
All nodes are persisted to the \texttt{mmr\_nodes} table in the same
database transaction as the event record. The MMR root is never stored
directly; it is recomputed from persisted nodes on demand, ensuring the
root always reflects actual persisted state.

\subsection{Epoch Chaining}

Events are assigned to epochs by monotonic per-organization sequence
number. When an epoch reaches \textsc{block\_size} events,
\texttt{\_finalize\_epoch()} is invoked within the same transaction:
(1) load all MMR nodes for the epoch; (2) compute MMR root by bagging
peaks; (3) serialize epoch metadata to canonical JSON; (4) compute epoch
hash; (5) persist the epoch record. The genesis epoch uses a fixed 32-byte
zero string as \texttt{prev\_epoch\_hash}.

\subsection{Two-Phase Commit Pipeline}

Agent decisions enter the ledger through a two-phase pipeline:

\textbf{Phase 1 --- Propose.} The agent submits a proposed action as a JSON
document. The pipeline serializes it, evaluates it against the policy
engine, and persists a \texttt{StagedEvent} record. If the proposal is
rejected and \texttt{EMIT\_BLOCKED\_ACTION} is enabled, a
\texttt{blocked\_action} audit event is immediately committed to the MMR.

\textbf{Phase 2 --- Commit.} An approved staged event is committed via
\texttt{\_insert\_event()}, which: assigns a monotonic sequence number;
serializes the payload; computes the leaf hash; appends to the MMR;
persists all new MMR nodes; persists the event record; and finalizes the
epoch if \textsc{block\_size} is reached. All operations occur within a
single database transaction using a single \texttt{db.commit()} at the
pipeline boundary.

\subsection{Verification}

An external auditor can verify any historical event independently:
(1)~\textit{Leaf verification}: recompute the leaf hash and compare to the
stored value; (2)~\textit{Inclusion proof}: verify the MMR proof path
against the epoch root; (3)~\textit{Epoch chain}: recompute each epoch
hash from a trusted anchor and compare to stored values.

%% ─────────────────────────────────────────────────────────────────────────────
\section{Implementation}

LICITRA-MMR is implemented in Python~3.11 using FastAPI, SQLAlchemy, and
PostgreSQL. The implementation is available at
\url{https://github.com/narendrakumarnutalapati/licitra-mmr-core} under MIT
license at tagged release v0.1.0-mvp.

\subsection{Key Design Decisions}

\textbf{MMR recomputed from persisted nodes.} The MMR root is never cached
in application memory. Every MMR operation loads the ordered node list from
the database, ensuring the root always reflects persisted state and any
direct database modification is immediately detectable.

\textbf{DEV\_MODE endpoint gating.} Tamper and reset endpoints return
HTTP~403 unless \texttt{DEV\_MODE=true} is explicitly set, with
\texttt{false} as the default. The \texttt{.env} file is excluded from
version control.

\textbf{Atomic transaction boundary.} The commit pipeline uses a single
\texttt{db.commit()} at the pipeline boundary. All intermediate writes use
\texttt{db.flush()} to assign database IDs without committing, ensuring any
failure results in a full rollback.

\textbf{SHA-256 via Python hashlib.} All cryptographic operations use
Python's \texttt{hashlib.sha256}, which wraps OpenSSL. SHA-256 is
FIPS~180-4 compliant and accepted by all major compliance frameworks,
making it preferable to BLAKE3 for compliance-oriented deployments despite
BLAKE3's superior throughput.

\subsection{Test Suite and Evidence Bundles}

The implementation includes 11 passing tests covering event proposal and
commitment, MMR node persistence, epoch finalization and hash chaining,
tamper detection, row deletion detection, and inclusion proof verification.
Five reproducible evidence scenarios are provided as JSON bundles:
S01 (normal commit), S02 (leaf tamper), S03 (epoch tamper), S04 (row
deletion), and S05 (cross-epoch verification).

%% ─────────────────────────────────────────────────────────────────────────────
\section{Competitive Analysis}

Table~\ref{tab:comparison} compares LICITRA-MMR against four related systems.

\begin{table*}[t]
\caption{Comparison of tamper-evident audit systems for agentic AI}
\label{tab:comparison}
\begin{tabular}{lcccc}
\toprule
Feature & LICITRA-MMR & immudb & CT Logs & Oktsec \\
\midrule
Data structure & MMR & Merkle tree & Merkle tree & Hash chain \\
Inclusion proof & $O(\log n)$ & $O(\log n)$ & $O(\log n)$ & None \\
Epoch anchoring & \checkmark & $\times$ & $\times$ & $\times$ \\
Per-org ledger & \checkmark & Partial & $\times$ & \checkmark \\
Pre-execution gate & $\times$ (SENTRY) & $\times$ & $\times$ & $\times$ \\
Explicit AI audit design & \checkmark & $\times$ & $\times$ & Partial \\
Canonical JSON spec & \checkmark versioned & $\times$ & $\times$ & $\times$ \\
Independent verification & \checkmark & \checkmark & \checkmark & $\times$ \\
License & MIT & Apache-2.0 & N/A & GPL-3.0 \\
Regulatory framing & EU AI Act, CO AI Act & General & PKI/TLS & None \\
\bottomrule
\end{tabular}
\end{table*}

\textbf{immudb}~\cite{immudb} is a general-purpose tamper-evident database
using a Merkle tree over a linear log. It provides inclusion proofs and
independent verification, making it the closest architectural peer. For use
cases requiring only tamper-evident storage without agentic governance
semantics, immudb is a mature alternative. LICITRA-MMR is distinguished by
its governance architecture: epoch anchoring, canonical JSON specification,
and integration with a pre-execution authorization layer.

\textbf{Certificate Transparency logs}~\cite{laurie2013} provide a globally
auditable append-only log for TLS certificates with gossip-based consistency
across independent operators. CT logs solve a different problem: proving
certificate existence in a publicly auditable multi-operator log. They
require external operators, have no per-organization namespacing, and are
not designed for high-frequency per-org event streams.

\textbf{Oktsec}~\cite{oktsec} provides agent identity verification and
message integrity using Ed25519 signatures and SHA-256 hash chains, covering
multiple OWASP Agentic Top~10 threat categories. LICITRA-MMR addresses a
complementary layer: epoch-anchored, independently verifiable commitment of
authorized agent decisions, rather than message-level integrity. The key
architectural distinction is that Oktsec's hash chain cannot produce
inclusion proofs without replaying the full chain; LICITRA's MMR enables
$O(\log n)$ verification of any single event.

%% ─────────────────────────────────────────────────────────────────────────────
\section{Theoretical Performance Evaluation}

All estimates assume Python~3.11 \texttt{hashlib.sha256} on an Intel
x86-64 processor (3.2~GHz, 32~GB RAM), with SHA-256 latency of
approximately 500~ns per call as measured via \texttt{timeit}
microbenchmark. Typical agent event payloads are assumed to be
approximately 200~bytes. Database I/O latency dominates the commit path;
cryptographic overhead is negligible in all scenarios.

\subsection{Inclusion Proof Path Length}

Table~\ref{tab:proof} shows proof path lengths at various epoch sizes. At
the default block size of 1,000, an inclusion proof requires 14 SHA-256
operations and occupies 462~bytes in binary encoding.

\begin{table}[ht]
\caption{MMR inclusion proof path length by epoch size}
\label{tab:proof}
\begin{tabular}{rrrr}
\toprule
Epoch size & Proof path & MMR nodes & Proof size \\
(events) & (hashes) & & (binary) \\
\midrule
100   & 8  & 197    & 264~B \\
500   & 13 & 994    & 429~B \\
1,000 & 14 & 1,994  & 462~B \\
2,000 & 15 & 3,994  & 495~B \\
10,000 & 17 & 19,995 & 561~B \\
\bottomrule
\end{tabular}
\end{table}

\subsection{Verification Cost at 1 Million Events}

With block size 1,000, one million events produce 1,000 epochs.
\textit{Epoch chain verification} requires 2 SHA-256 operations per epoch
(2,000 total), completing in approximately 1~ms. \textit{Per-event inclusion
proof} requires 14 SHA-256 operations and completes in approximately 7~$\mu$s
per event. \textit{Full leaf rehash} of all 1 million events requires
approximately 2~seconds at 2~$\mu$s per event. \textit{Sample verification}
of 100 randomly selected events requires approximately 0.9~ms.

\subsection{Commit Overhead}

The amortized cryptographic overhead per commit is approximately 6 SHA-256
operations ($\approx3~\mu$s). The dominant commit cost is database I/O:
three write operations per commit.

\subsection{Storage at Scale}

At 1 million events with block size 1,000, the estimated storage footprint
is approximately 564~MB: MMR nodes~$\approx$217~MB (1,994,000 node rows),
event rows~$\approx$347~MB (1,000,000 rows at~$\approx$364~bytes each),
epoch rows~$<1$~MB. Storage scales linearly with event count. The MMR node
overhead is approximately $2\times$ the leaf count, consistent with the
asymptotic minimum for append-only binary hash tree constructions.

%% ─────────────────────────────────────────────────────────────────────────────
\section{Architectural Non-Goals}

The following properties are explicitly outside the scope of LICITRA-MMR.

\textbf{Not a consensus system.} There is one writer and one ledger per
organization. Consistency is enforced by database transaction semantics,
not a consensus protocol.

\textbf{Not a blockchain.} The epoch chain is structurally similar to a
blockchain header chain but is not a blockchain. There is no proof of work,
proof of stake, token, peer-to-peer network, or distributed ledger.

\textbf{Not Byzantine fault tolerant.} The system assumes the database
storage layer faithfully preserves byte-level fidelity of stored values.
It does not protect against a Byzantine storage layer.

\textbf{Not preventing application-layer falsification.} LICITRA-MMR
commits whatever the application layer submits. Pre-execution integrity is
the scope of LICITRA-SENTRY, not LICITRA-MMR.

\textbf{Not availability-protecting.} Availability is the responsibility of
the deployment infrastructure.

%% ─────────────────────────────────────────────────────────────────────────────
\section{Limitations and Future Work}
\label{sec:limitations}

\subsection{Float Normalization Gap}

The canonical JSON specification does not normalize floating-point values.
The integer \texttt{1} and the float \texttt{1.0} produce different byte
sequences and different leaf hashes. For the primary use case of agent
decision logging, this is unlikely to cause verification failures in
practice. Full float normalization per RFC~8785~\S3.2.2 is planned for
CANONICAL\_JSON\_SPEC~v1.1.

\subsection{Unsigned Epoch Roots}

The current implementation computes epoch hashes but does not sign them
with a private key. An attacker who controls both the database and the
application server can in principle reconstruct a consistent but false
history. Ed25519 epoch signing is planned for v0.2. The epoch record
schema reserves a \texttt{signature} field for this purpose.

\subsection{Single-Operator Trust Model}

LICITRA-MMR is designed for single-operator deployments. Multi-party
witnessing --- where multiple independent operators each hold a copy of the
current epoch hash --- is a direction for future work, analogous to the
gossip protocol in Certificate Transparency~\cite{laurie2013}.

\subsection{Pre-Execution Integrity}

LICITRA-MMR records what was committed. It does not guarantee that what was
committed accurately reflects what the agent was authorized to do.
Pre-execution integrity is the scope of LICITRA-SENTRY.

%% ─────────────────────────────────────────────────────────────────────────────
\section{Related Work}

\textbf{Merkle Mountain Ranges.} MMRs were introduced by Peter
Todd~\cite{todd2013} and have been adopted in blockchain systems including
Grin and Beam. The formal optimality of the MMR for append-only
accumulators was established by Bonneau et al.~\cite{bonneau2025}.
LICITRA-MMR applies the MMR structure to per-organization agentic AI audit
logs, a context not previously addressed in the literature.

\textbf{Certificate Transparency.} Laurie et al.~\cite{laurie2013}
introduced CT logs for publicly auditing TLS certificate issuance.
LICITRA-MMR shares the append-only authenticated structure but differs in
trust model, scope, and operational requirements.

\textbf{immudb.} Codenotary's immudb~\cite{immudb} provides a
tamper-evident database using a Merkle tree over an append-only log. It is
the closest architectural peer; distinctions are described in Section~6.

\textbf{Oktsec.} The Oktsec framework~\cite{oktsec} provides agent identity
verification and message integrity. LICITRA-MMR addresses a complementary
layer at the authorization and audit level.

\textbf{NIST AI Risk Management Framework.} The NIST AI RMF~\cite{nistai}
provides governance guidance including auditability requirements.
LICITRA-MMR provides a technical implementation of that requirement.

%% ─────────────────────────────────────────────────────────────────────────────
\section{Conclusion}
\label{sec:conclusion}

We have presented LICITRA-MMR, an open-source cryptographic ledger primitive
for agentic AI runtime accountability. The system combines a Merkle Mountain
Range with per-organization epoch anchoring, a versioned canonical JSON
specification, and an atomic two-phase commit pipeline to provide forensic
integrity guarantees that mutable logging infrastructure cannot offer.

The three core contributions --- the canonical specification, the
MMR-based epoch-anchored ledger, and the atomic commit pipeline --- address
the gap between what emerging AI regulation implies (post-incident
reconstructability with sufficient integrity for audit) and what most
production agentic AI stacks currently provide (operator-controlled mutable
logs).

LICITRA-MMR is the foundational layer of the broader LICITRA cryptographic
governance stack. The next layer, LICITRA-SENTRY, adds a zero-trust control
plane with pre-execution semantic contracts, agent identity verification,
content inspection, and Chain of Intent authorization.

\subsection{Deployment Considerations}

Organizations deploying LICITRA-MMR should consider: (1)~\textit{Epoch hash
anchoring} --- store at least one epoch hash off-site via a notarization
service, timestamping service, or publication to a public ledger, to
prevent an operator from reconstructing a consistent false history; (2)~%
\textit{Key management preparation} --- establish KMS or HSM infrastructure
in advance of the v0.2 signing upgrade; (3)~\textit{Backup and retention}
--- the MMR node table grows at approximately $2\times$ the event count;
(4)~\textit{Regulatory export} --- use the JSON evidence bundle format from
the \texttt{licitra-mmr-evidence} repository for regulatory submissions.

\textbf{Reproducibility.} All code, test suites, and evidence bundles are
available at \url{https://github.com/narendrakumarnutalapati/licitra-mmr-core}.
The tagged release v0.1.0-mvp and all evidence bundles are permanently
archived at \url{https://doi.org/10.5281/zenodo.18796529}.

%% ─────────────────────────────────────────────────────────────────────────────
\bibliographystyle{ACM-Reference-Format}
\begin{thebibliography}{9}

\bibitem{todd2013}
P.~Todd.
\newblock Merkle Mountain Ranges.
\newblock OpenTimestamps, 2013.
\newblock \url{https://github.com/opentimestamps/opentimestamps-server/blob/master/doc/merkle-mountain-range.md}

\bibitem{laurie2013}
B.~Laurie, A.~Langley, and E.~Kasper.
\newblock Certificate Transparency.
\newblock RFC~6962, Internet Engineering Task Force, 2013.
\newblock \url{https://www.rfc-editor.org/rfc/rfc6962}

\bibitem{bonneau2025}
J.~Bonneau, J.~Chen, M.~Christ, and I.~Karantaidou.
\newblock Merkle Mountain Ranges are Optimal: On Witness Update Frequency for
  Cryptographic Accumulators.
\newblock In \textit{Advances in Cryptology -- CRYPTO~2025}, LNCS vol.~16001,
  Springer, 2025.
\newblock \url{https://eprint.iacr.org/2025/234}

\bibitem{immudb}
Codenotary.
\newblock immudb: Immutable Database.
\newblock \url{https://immudb.io}

\bibitem{oktsec}
Oktsec.
\newblock Agent Security Framework (v0.1).
\newblock \url{https://github.com/oktsec/oktsec}
\newblock Accessed: February 2026.

\bibitem{nistai}
National Institute of Standards and Technology.
\newblock Artificial Intelligence Risk Management Framework (AI~RMF~1.0).
\newblock \url{https://doi.org/10.6028/NIST.AI.100-1}, 2023.

\bibitem{euaiact}
European Parliament.
\newblock Regulation (EU) 2024/1689 on Artificial Intelligence (EU AI Act).
\newblock Official Journal of the European Union, 2024.

\bibitem{coloradoai}
Colorado General Assembly.
\newblock SB~24-205: Colorado Artificial Intelligence Act.
\newblock Effective June~30, 2026.
\newblock \url{https://leg.colorado.gov/bills/sb24-205}

\end{thebibliography}

\end{document}
